\documentclass[10pt,a4paper]{article}
% ---------- common preamble : CM3191 organic documents (XeLaTeX) ----------
\usepackage[a4paper,top=16mm,bottom=16mm,left=16mm,right=16mm,headheight=14pt]{geometry}
\usepackage{fontspec}
\usepackage{xeCJK}
\setmainfont{TeX Gyre Termes}
\setsansfont{TeX Gyre Heros}
\setCJKmainfont{Noto Serif CJK KR}
\setCJKsansfont{Noto Sans CJK KR}
\xeCJKsetup{CJKspace=true, CJKecglue={}}
\usepackage{amsmath,amssymb}
\usepackage[version=4]{mhchem}
\usepackage{chemfig}
\usepackage{siunitx}
\sisetup{detect-all, per-mode=symbol, range-phrase=--, range-units=single}
\usepackage{tikz}
\usetikzlibrary{arrows.meta,calc,positioning,shapes.geometric,decorations.pathmorphing,patterns}
\usepackage{tabularx,array,multirow,booktabs}
\usepackage{enumitem}
\usepackage{titlesec}
\usepackage{fancyhdr}
\usepackage{xcolor}
\usepackage{pifont}
\usepackage[hidelinks]{hyperref}
\setlength{\parindent}{0pt}
\setlength{\parskip}{3pt}
\linespread{1.08}

\newcommand{\doctag}{}
\pagestyle{fancy}\fancyhf{}
\lhead{\footnotesize CM3191 Organic \textemdash\ \doctag}
\rhead{\footnotesize KIM DONGHWAN \textbullet\ B7 \textbullet\ Thu}
\cfoot{\footnotesize \thepage}
\renewcommand{\headrulewidth}{0.4pt}

\titleformat{\section}{\large\bfseries}{\thesection}{0.6em}{}
\titlespacing*{\section}{0pt}{9pt}{4pt}
\titleformat{\subsection}{\normalsize\bfseries}{\thesubsection}{0.5em}{}
\titlespacing*{\subsection}{0pt}{7pt}{3pt}

% ---------- form primitives ----------
\newcommand{\wl}[1][4cm]{\rule[-0.6ex]{#1}{0.4pt}}
\newcommand{\cbox}{$\square$}
\newcommand{\tr}{\rule[-3mm]{0pt}{9mm}}
\newcommand{\trs}{\rule[-2mm]{0pt}{7mm}}
\newcommand{\ruled}[1]{%
  \par\vspace{3mm}%
  \foreach \i in {1,...,#1}{\noindent\rule{\linewidth}{0.35pt}\par\vspace{7mm}}%
  \vspace{-7mm}}
\newcommand{\confirm}[1]{\textbf{[CONFIRM: #1]}}

% ---------- boxed callout ----------
\newenvironment{callout}[1]
  {\par\vspace{2mm}\noindent\begin{tabular}{|@{\hspace{2.5mm}}p{\dimexpr\linewidth-7.4mm\relax}@{\hspace{2.5mm}}|}\hline
   \rule{0pt}{4.5mm}\textbf{#1}\par\vspace{0.5mm}\small}
  {\par\vspace{1.5mm}\\\hline\end{tabular}\par\vspace{2mm}}

% ---------- title block ----------
\newcommand{\manualtitle}[3]{%
\begin{center}
{\Large\bfseries #1}\\[1.5mm]
{\large #2}\\[2mm]
{\footnotesize #3}
\end{center}
\vspace{-1mm}\rule{\linewidth}{1pt}\vspace{1mm}
}

% ---------- GTA sign-off ----------
\newcommand{\gtasign}{%
\par\vspace{3mm}
\renewcommand{\arraystretch}{1.4}
\noindent\begin{tabularx}{\linewidth}{|X|X|X|}
\hline
\multicolumn{3}{|l|}{\textbf{To be completed at the end of the session}}\\ \hline
Student signature & GTA name & GTA signature \\ \hline
\tr & & \\ \hline
\end{tabularx}
\renewcommand{\arraystretch}{1}
\par\vspace{1mm}
{\footnotesize This is a working log. The \textbf{official} datasheet (lab manual p.\,7) must still be printed, signed by the GTA and submitted with the report.}
}

% ---------- chemfig defaults ----------
\setchemfig{atom sep=1.9em, bond style={line width=0.6pt}}

\renewcommand{\doctag}{Organic Part Roadmap}
\setlength{\parskip}{4pt}
\linespread{1.12}

\begin{document}

\manualtitle{CM3191 유기 파트 로드맵}
{Week 7--13 일정, 실험별 준비, 확인할 질문}
{근거: CM3191 Organic Lecture and Lab schedule, Schlenk line 조 배정 xlsx(Thursday Group A, B7 KIM DONGHWAN), Organic Lab Manual AY26/27, NUS 학사력 AY2026/27. 평가 방식은 schedule에 ``7주차 전에 업데이트''로만 되어 있어 아직 반영하지 않았습니다.}

\section*{1. 주간 일정}
\renewcommand{\arraystretch}{1.3}
\begin{tabularx}{\linewidth}{|c|>{\footnotesize}p{2.1cm}|>{\footnotesize}X|>{\footnotesize}X|>{\footnotesize}p{3.6cm}|}
\hline
\textbf{주} & \textbf{날짜} & \textbf{강의 (화 10--12)} & \textbf{실험 (목, Dr Hoang)} & \textbf{미리 할 일} \\ \hline
7 & 화 9/29 \newline 목 10/1 &
\textbf{Schlenk line 실습} 10:50--11:35, GTL (본인 조) &
\textbf{Exp.\,1 Suzuki} &
Schlenk 준비 문서, Suzuki 3종 문서, 공식 datasheet(p.\,7) 인쇄 \\ \hline
8 & 화 10/6 \newline 목 10/8 &
Schlenk 실습은 Group B 차례 (본인 해당 없음) &
\textbf{Exp.\,2 aldol--Diels--Alder} &
매뉴얼 pp.\,8--13, L2. datasheet(p.\,13) 인쇄. TLC 그림 공간 준비 \\ \hline
9 & 화 10/13 \newline 목 10/15 &
입체선택적 환원 실험 강의 (UTown) &
\textbf{Exp.\,4 Grignard 1} (Schlenk line, 반응·workup·TLC) &
pp.\,19--24, Schlenk 실습 메모. datasheet(p.\,27) 인쇄 \\ \hline
10 & 화 10/20 \newline 목 10/22 &
Flow chemistry 시연 (GTL) &
\textbf{Exp.\,3 환원} (NaBH\textsubscript{4}) &
pp.\,14--18. datasheet(p.\,18) 인쇄 \\ \hline
11 & 목 10/29 & --- & \textbf{Grignard 2} (컬럼 크로마토그래피) & pp.\,24--26, 9주차 TLC 결과 지참 \\ \hline
12 & 목 11/5 & --- & Make-up lab & 필요 시 \\ \hline
13 & 화 11/10 (추정) & \textbf{Written test} & --- & 월요일 11/9는 Deepavali 대체공휴일 \\ \hline
\end{tabularx}
\renewcommand{\arraystretch}{1}
{\footnotesize 날짜는 NUS 학사력(7주차 = 9/28--10/2)으로 계산했습니다. 9--10주차 순서는 Group A 기준입니다(Grignard 먼저, 환원 나중). 금 10/9는 NUS Well-Being Day로 목요일 실험에는 영향이 없습니다.}

\section*{2. 실험별 한눈에 보기}
\renewcommand{\arraystretch}{1.3}
\begin{tabularx}{\linewidth}{|>{\footnotesize}p{2.4cm}|>{\footnotesize}X|>{\footnotesize}X|>{\footnotesize}p{3.0cm}|}
\hline
\textbf{실험} & \textbf{핵심 반응과 양} & \textbf{핵심 기술} & \textbf{분석} \\ \hline
Exp.\,1 Suzuki &
4-bromobenzoic acid 1.69 mmol + \ce{PhB(OH)2} 1.21 eq, Pd 0.015 mol\%, 물, 70\,\si{\celsius} 30 min. 이론 0.335 g &
pH로 녹이고 꺼내기, 재결정(EtOH) &
IR, \textsuperscript{1}H NMR \\ \hline
Exp.\,2 aldol--DA &
benzil + 1,3-diphenylacetone 각 1.00 mmol \textrightarrow\ tetracyclone(보라) \textrightarrow\ styrene 3.5 eq. 이론 0.488 g (3A/3B) 또는 0.460 g (4) &
한 용기 두 반응, TLC 두 번, 혼합 용매 재결정 &
녹는점, IR, \textsuperscript{1}H NMR, X선 구조 해석 \\ \hline
Exp.\,4 Grignard &
6-methyl-5-hepten-2-one 2.98 mmol + vinylMgBr 2.5 eq, THF. 이론 linalool 0.460 g &
Schlenk line, 주사기 이동, 추출, 컬럼 &
TLC(KMnO\textsubscript{4}), IR, \textsuperscript{1}H NMR \\ \hline
Exp.\,3 환원 &
4-\textit{tert}-butylcyclohexanone 3.24 mmol + \ce{NaBH4} 3.3 eq(hydride 기준 약 13 eq), EtOH. 이론 0.507 g &
온도 조절, \ce{H2} 발생 quench, 추출, 회전증발 &
TLC(anisaldehyde), IR, \textsuperscript{1}H NMR(cis/trans 비) \\ \hline
\end{tabularx}
\renewcommand{\arraystretch}{1}
{\footnotesize 수치는 매뉴얼 표의 양으로 직접 계산했습니다(이전 정리와 동일).}

\section*{3. 반복해서 나오는 기술 (한 번 익히면 계속 씀)}
\begin{itemize}[leftmargin=5mm,itemsep=0.5mm]
\item \textbf{재결정}: Exp.\,1(에탄올 단일 용매), Exp.\,2(에탄올/에틸아세테이트 혼합). ``최소량의 뜨거운 용매''라는 원리는 같습니다.
\item \textbf{TLC}: Exp.\,2, 3, 4. 각 datasheet에 판 그림과 거리(cm)를 적어야 하므로 \textbf{자와 연필}을 매번 가져가세요.
\item \textbf{추출과 회전증발}: Exp.\,3, 4. 분별 깔때기 사용법은 Exp.\,3(10주차)보다 Exp.\,4(9주차)에서 먼저 나옵니다.
\item \textbf{\textsuperscript{1}H NMR 해석}: 네 실험 모두. 적분과 갈라짐 읽는 법은 ``기초부터 해설'' 8장에 있습니다.
\end{itemize}

\section*{4. 모아 둔 확인 질문 (TA 또는 강사)}
\begin{enumerate}[leftmargin=6mm,itemsep=0.5mm]
\item 유기 실험도 risk assessment 제출이 필요한가? 필요하면 양식과 마감은?
\item 보고서·proforma 형식과 마감 (평가 방식 업데이트 확인)
\item 유기 실험실 장소 (schedule에는 강의 장소만 있음)
\item Exp.\,1 NMR 용매가 DMSO-$d_6$인가? (매뉴얼은 ``the NMR solvent''로만 표기)
\item 유기 실험실의 폐액 분류: 금속 함유 수용액, 할로젠/비할로젠 유기 폐액
\item Exp.\,2 Safety에 적힌 iodine과 methylene chloride는 어디에 쓰는가? (절차에는 없음)
\item Schlenk line: 코크 방향, 불활성 기체 종류, cold trap 냉각 방식
\item Grignard 조생성물을 9주차부터 11주차까지 어떻게 보관하는가?
\item 조 배정표의 2--3명 묶음(본인은 B7, D1 Koh Liang Han과 한 묶음)이 Schlenk line 공유 단위인가? Grignard 반응은 각자 하는가, 함께 하는가?
\end{enumerate}

\section*{5. 준비된 문서 (드롭박스 \texttt{/chem/CM3191/organic/})}
\renewcommand{\arraystretch}{1.2}
\begin{tabularx}{\linewidth}{|l|X|}
\hline
\texttt{1/} Suzuki Bench Sheets & 체크리스트, 영문·국문 매뉴얼, 작업용 datasheet, 장치 그림 \\ \hline
\texttt{1/} Suzuki Concepts & 촉매 순환, 염기, 물 용매, TON/TOF, 후처리, 스펙트럼, 연습문제 \\ \hline
\texttt{1/} Suzuki Beginner Guide & 구조식 읽기부터 시작하는 입문 해설 \\ \hline
\texttt{1/} Suzuki Report Draft & 보고서 초안 (결과 자리는 회색 상자) \\ \hline
\texttt{1/} Schlenk Line Prep & 9/29 실습 준비 (Grignard 예습 겸용) \\ \hline
\texttt{2/}, \texttt{3/}, \texttt{4/} Report Draft & aldol--DA, 환원, Grignard 보고서 초안 (\texttt{report\_common.tex} 필요) \\ \hline
이 로드맵, \texttt{README.md} & 전체 일정, 질문 목록, 폴더 안내 \\ \hline
\end{tabularx}
\renewcommand{\arraystretch}{1}

\end{document}
